{{دیوان
| عنوان = ابوالعلامعرّیؔ
| ماخذ = ویکی ماخذ
| ماخذ_لنک = https://ur.wikisource.org/wiki/%D8%A7%D8%A8%D9%88%D8%A7%D9%84%D8%B9%D9%84%D8%A7%D9%85%D8%B9%D8%B1%D9%91%DB%8C%D8%94
| مقابلہ = اقبال سائبر لائبریری (اقبال اکادمی پاکستان)
| مقابلہ_لنک = https://www.iqbalcyberlibrary.net
}}
<poem>
کہتے ہیں کبھی گوشت نہ کھاتا تھا معّری
پھل پھُول پہ کرتا تھا ہمیشہ گزر اوقات

اک دوست نے بھُونا ہوا تِیتر اسے بھیجا
شاید کہ وہ شاطِر اسی ترکیب سے ہو مات

یہ خوانِ تر و تازہ معّری نے جو دیکھا
کہنے لگا وہ صاحبِ عفران و لزومات

اے مُرغکِ بیچارہ! ذرا یہ تو بتا تُو
تیرا وہ گُنہ کیا تھا یہ ہے جس کی مکافات؟

افسوس، صد افسوس کہ شاہیں نہ بنا تُو
دیکھے نہ تری آنکھ نے فطرت کے اشارات

تقدیر کے قاضی کا یہ فتویٰ ہے ازل سے
ہے جُرمِ ضعیفی کی سزا مرگِ مفاجات!
</poem>

٭ عفران۔ رسالتہ الغفران، معری کی ایک مشہور کتاب کا نام ہے
